\scanpage{17}\mspage{16}
% ink heading and first remark
\uline{Remarques} \textcircled{1} \deleted{\unsure{Considér|Considé}} sur la notion de tronçon. On cherche une sous-cat. pleine \inserted{$\unsure{\mathscr T|\mathcal T}$} de celle des ens. ordonnés, satisfaisant

% ink item a, two lines and qualification at the right
a) \deleted{\unsure{S}}Tout $\mathcal L\in\mathrm{Ob}\,\mathscr T$ est filtrant croissant et filtrant décroissant, et \deleted{non vide} \inserted{de cardinal $\ne1$} (donc $\geq2$)

% ink item b
b) Si $x,y\in\deleted{\unsure{T}}\mathcal L_{\deleted{\unsure{,|g}}}\in\deleted{\unsure{O}}\mathrm{Ob}\,\mathscr T$, et $x<y$, $\exists$
\[
 z\in L,\qquad x<z<y\qquad\text{(divisibilité)}
\]

% ink item c
c) Si $x\in L\in\unsure{\mathrm{Ob}|\mathrm D}\,\mathscr T$, \deleted{alors} et si
\[
 L_{>x}\ne\varnothing,\quad\text{alors}\quad L_{>x}\in\mathrm{Ob}\,\mathscr T\,.
\]

% ink item c prime
c$'$) Dual de c)

% ink middle and lower third, Il y a
Il y a une plus grande $\mathscr T$ qui satisfait ces conditions, et $\mathscr T$ est stable par passage à l'ordre opposé. Les objets de $\mathscr T$ sont les tronçons \unsure{ordonnés}.

% ink bottom, second remark
\textcircled{2} Il faudrait expliciter \unsure{quelle|quelles} \unsure{condition|conditions} la relation $R(a,b,c)$ doit satisfaire, pour définir une structure de tronçon \unsure{biordonné|biordonnien}\unsure{,} et itou pour la donnée de $\unsure{\mathscr S|S}_0\subset\unsure{\mathrm P|\mathscr P}_2(L)$ et les $R_{\unsure{E|\mathcal E}}$\unsure{,} ($E\in\mathscr S_0$) sur les $C(E)$ \ldots
